\documentclass[11pt,a4paper]{article}
\usepackage[utf8]{inputenc}
\usepackage[T1]{fontenc}
\usepackage{geometry}
\usepackage{amsmath,amssymb}
\usepackage{booktabs}
\usepackage{graphicx}
\usepackage[numbers]{natbib}
\usepackage{listings}
\usepackage[colorlinks=true,linkcolor=blue,citecolor=blue,urlcolor=blue]{hyperref}

\geometry{
  left=2.5cm,
  right=2.5cm,
  top=2.5cm,
  bottom=2.5cm
}

\title{\textbf{NRR-Boundary: Same Operators, Opposite Effects by Provider and Order}}

\author{
  Kei Saito\thanks{ORCID: 0009-0006-4715-9176} \\
  Independent Researcher, Japan \\
  \texttt{kei.saito.research@gmail.com}
}

\date{March 2026 \\[0.5em]
\textit{Part of the Non-Resolution Reasoning (NRR) research program.}}

\begin{document}
\maketitle

\begin{center}
\textcopyright\ 2026 Kei Saito.
Licensed under CC BY 4.0.\\
\url{https://creativecommons.org/licenses/by/4.0/}
\end{center}

\begin{abstract}
NRR operators are often discussed as if their quality effects were intrinsic to each operator. This paper evaluates that assumption under fixed paired comparisons and reports a boundary map instead of a single winner ranking. Under tested conditions (spec-to-code tasks, complexity c=0.6, temperatures t=0.0 and t=0.3 for Stage B, and t=0.0 for ordered-combination follow-up), operator effects vary by provider and can reverse by application order. Stage B (single operators) shows provider-dependent sign differences, including strong degradation for Gemini-lazy and positive improvement means for conditional across all tested providers. Ordered-combination follow-up (reps=3) shows that the same operator set can move in opposite directions when order changes, with the largest swing observed in Gemini rework cost for conditional\_branch vs branch\_conditional. We present provider-separated and provider-balanced outputs, explicit sign-flip boundaries, and reproducible run artifacts. These findings imply that operator deployment recommendations should be qualified by provider and order context.
\end{abstract}

\noindent\textbf{Keywords:} Non-Resolution Reasoning, operator interaction, sign flip boundary, provider heterogeneity, ordered composition

\section{Introduction}
\label{sec:intro}

Prior NRR papers defined delayed commitment and its operational implications~\cite{saito2025nrr,saito2026phi,saito2026ime}. This manuscript asks a narrower engineering question: under fixed conditions, where do NRR operators improve quality, where do they degrade it, and where do they become non-effective?

This question matters because operator recommendations are increasingly reused across models as if they were portable defaults. In real deployment, however, teams often switch provider endpoints or compose multiple operators in pipelines. If effect direction changes with provider or with operator order, a seemingly benign prompt policy can become a quality or cost risk. The practical target of this paper is therefore not a universal operator ranking, but a deployable boundary map that answers: ``safe where, risky where, and under which configuration?''

The same issue also affects empirical reliability. Without explicit provider/order boundaries, replication failures are easy to misclassify as random noise, while they may actually be deterministic configuration effects. A boundary map is therefore useful for both operations and science: operations teams can avoid known high-risk cells before rollout, and researchers can separate ``true disagreement'' from ``different cell, different behavior'' when comparing results across studies.

In this paper, ``NRR operators'' are user-side update policies (e.g., conditional, hierarchical, delta, branch, lazy) that delay or structure commitment during iterative task updates. Delayed commitment means not forcing an early single interpretation when evidence is still unstable. MST evaluates whether these operators change downstream quality signals under identical baseline-vs-pattern conditions. In filenames and run identifiers, ``MST'' is retained as the internal experiment code name for reproducibility tracking, while the paper-facing name is NRR-Boundary.

NRR is not an anti-LLM framework.
NRR does not replace standard LLM use.
NRR optimizes when to commit and when to defer, under explicit conditions.

We therefore avoid global claims such as ``operator X is best'' and instead report a boundary map over provider, pattern, and order.

\textbf{Series Consistency Statement:}
Across the NRR series, unsupported interpretations are not injected at a fixed evidence state, hypothesis-set expansion is evidence-gated, and unobserved possibilities remain explicitly unresolved. In MST, this policy is operationalized through fixed paired comparisons and explicit reporting of improvements, regressions, and sign-flip boundaries.

\subsection{Contributions}

\begin{enumerate}
  \item A fixed paired protocol for measuring operator effects under identical scenario conditions.
  \item A confirmatory single-operator map (Stage B) with provider-separated and provider-balanced outputs.
  \item An ordered-combination follow-up (reps=3) that quantifies provider-by-order interactions.
  \item Boundary reporting with sign-flip counts and non-effective regions instead of winner-only summaries.
\end{enumerate}

\section{Related Work}
\label{sec:related}

Recent LLM evaluation work emphasizes benchmark breadth and standardized reporting (e.g., HELM)~\cite{liang2022helm}. Interactive quality benchmarks such as MT-Bench and Chatbot Arena also show that model ranking depends on setup details, including prompt/judge protocol and comparison context~\cite{zheng2023mtbench}. Open instruction-following evaluations (e.g., AlpacaEval) similarly demonstrate that preference-oriented outcomes vary with evaluator design~\cite{dubois2024alpacaeval}.

Prompt-engineering literature has repeatedly shown sensitivity to formatting, reasoning scaffolds, and decoding-time strategy. Chain-of-thought prompting~\cite{wei2022cot}, zero-shot reasoning triggers~\cite{kojima2022zeroshot}, and self-consistency decoding~\cite{wang2023selfconsistency} can each change output quality materially. Calibration studies also report that small prompt-template changes alter few-shot behavior~\cite{zhao2021calibrate}. These findings motivate a boundary-first evaluation style: when prompt-side controls are powerful, their effects should be measured under fixed and auditable protocols rather than assumed to transfer unchanged.

This paper is complementary to those lines. Rather than proposing a new global benchmark, we fix a paired protocol and map where operator-induced quality changes are positive, negative, or non-effective. In this sense, the target output is a boundary map over provider, pattern, and order under frozen run conditions.

\section{Method}
\label{sec:method}

\subsection{Paired Evaluation Policy}

Each comparison uses the same scenario and scoring schema for baseline vs pattern arms. Outcome-based retry is disabled (\texttt{max\_quality\_retries=0}). We report both gains and regressions.

Primary confirmatory metrics are:
\begin{itemize}
  \item \texttt{misinterpretation\_rate}
  \item \texttt{trace\_completeness}
  \item \texttt{update\_consistency}
  \item \texttt{internal\_coherence\_error\_rate}
  \item \texttt{self\_contradiction\_rate}
  \item \texttt{rework\_cost\_tokens}
\end{itemize}

IRP-related diagnostics (\texttt{stability\_variance}, \texttt{irp\_monotonic\_progress\_rate}, \texttt{irp\_oscillation\_rate}) are retained as secondary diagnostics.

\subsection{Stage B (Primary Confirmatory)}

Stage B run directory:
\texttt{results/main\_c06\_5patterns\_3providers\_20260302T054857Z}

Configuration summary:
\begin{itemize}
  \item patterns: conditional, hierarchical, delta, branch, lazy
  \item providers: openai, gemini, anthropic
  \item temperatures: t=0.0 and t=0.3
  \item reps: 3 per scenario cell
\end{itemize}

Observed scale:
\begin{itemize}
  \item 30 raw files (10 per provider)
  \item 270 paired comparisons
  \item 2430 pair-metric rows
\end{itemize}

Confirmatory multiplicity control follows fixed Holm family size $n=30$ (5 patterns $\times$ 6 confirmatory metrics).

\subsection{Ordered-Combination Follow-up (Secondary)}

Follow-up run directory:
\texttt{results/combo\_probe\_c06\_ordered8\_rep3\_20260303T020142Z}

Ordered pairs:
\begin{itemize}
  \item conditional\_delta / delta\_conditional
  \item conditional\_hierarchical / hierarchical\_conditional
  \item conditional\_branch / branch\_conditional
  \item delta\_hierarchical / hierarchical\_delta
\end{itemize}

Configuration summary:
\begin{itemize}
  \item providers: openai, gemini, anthropic
  \item temperature: t=0.0 only
  \item reps: 3
\end{itemize}

Observed scale:
\begin{itemize}
  \item 24 raw files (8 patterns $\times$ 3 providers)
  \item 216 paired comparisons
  \item 1944 pair-metric rows
\end{itemize}

This follow-up is treated as secondary boundary mapping, and effect-size/sign-flip reporting is prioritized over universalized claims.

\subsection{Statistical Positioning}

Stage B is treated as the confirmatory test family with fixed Holm correction size $n=30$. The ordered-combination follow-up is reported as descriptive boundary mapping. Because the follow-up design was selected after Stage B observations, it is not treated as a new confirmatory family in this manuscript.

For Stage B confirmatory tests, we apply two-sided Wilcoxon signed-rank tests to pair-level improvement values (provider-pooled within each pattern-metric cell), followed by Holm-Bonferroni correction with family size $n=30$.

\section{Results I: Single-Operator Effects (Stage B)}
\label{sec:results_stageb}

\subsection{Gate Check}

Using all pair-metric rows in the Stage B full run:
\begin{itemize}
  \item $\tau$ trace pass rate: 0.988889
  \item parse-fail session rate: 0.005556
  \item non-effective rate: 0.722222
\end{itemize}

Under tested thresholds ($\tau\geq0.95$, parse-fail $<0.10$, non-effective $<1.0$), Stage B passes.

\subsection{Provider-by-Pattern Summary (Key Metrics)}

Table~\ref{tab:stageb_key} reports improvement means for two high-discriminability metrics (higher is better; negative means degradation).

\begin{table}[h]
\centering
\small
\begin{tabular}{llll}
\toprule
Provider & Pattern & Misinterpretation improvement & Rework improvement (tokens) \\
\midrule
OpenAI & conditional & +0.087963 & +82.166667 \\
OpenAI & hierarchical & +0.078704 & +80.500000 \\
OpenAI & delta & +0.041667 & +48.500000 \\
OpenAI & branch & +0.027778 & +35.111111 \\
OpenAI & lazy & +0.006944 & -0.444444 \\
\midrule
Gemini & conditional & +0.118519 & +221.666667 \\
Gemini & hierarchical & -0.075926 & -194.611111 \\
Gemini & delta & +0.022222 & -21.944444 \\
Gemini & branch & -0.035185 & -121.444444 \\
Gemini & lazy & -0.412963 & -780.000000 \\
\midrule
Anthropic & conditional & +0.027778 & +64.111111 \\
Anthropic & hierarchical & +0.033333 & +77.555556 \\
Anthropic & delta & -0.055556 & -106.444444 \\
Anthropic & branch & +0.005556 & +12.166667 \\
Anthropic & lazy & +0.033333 & +77.333333 \\
\bottomrule
\end{tabular}
\caption{Stage B provider-separated improvement means for key metrics.}
\label{tab:stageb_key}
\end{table}

Key observation under tested conditions: conditional showed positive improvement means on these key metrics across all tested providers, while other patterns show provider-dependent reversals.

\begin{table}[h]
\centering
\small
\begin{tabular}{lrr}
\toprule
Pattern & Misinterpretation improvement (balanced) & Rework improvement (balanced, tokens) \\
\midrule
conditional & +0.0781 & +122.6481 \\
hierarchical & +0.0120 & -12.1852 \\
delta & +0.0028 & -26.6296 \\
branch & -0.0006 & -24.7222 \\
lazy & -0.1242 & -234.3704 \\
\bottomrule
\end{tabular}
\caption{Stage B provider-balanced improvement means for key metrics (equal provider weighting).}
\label{tab:stageb_key_balanced}
\end{table}

Table~\ref{tab:stageb_key_balanced} makes the provider-balanced view explicit. Because Stage B uses equal provider counts (18/provider), provider-balanced key-metric means align closely with the pooled Stage B summaries used in confirmatory testing.

\begin{figure}[h]
\centering
\includegraphics[width=\linewidth]{stageb_provider_separated_heatmaps_v9_readable.png}
\caption{Stage B provider-separated heatmaps across five single operators (processed/all). Each metric panel uses its own symmetric color scale centered at zero; color intensity is comparable within a panel (not across different metrics).}
\label{fig:stageb_provider_heatmap}
\end{figure}

\begin{figure}[h]
\centering
\includegraphics[width=\linewidth]{stageb_sign_flip_boundaries.png}
\caption{Stage B sign-flip boundary map (processed/all).}
\label{fig:stageb_signflip}
\end{figure}

Stage B sign-flip boundaries appear in 20/45 pattern-metric cells, concentrated in branch, delta, hierarchical, and lazy patterns. Conditional does not show sign-flip boundaries in Stage B under this setup. This directly supports the transition from single-operator mapping in Fig.~\ref{fig:stageb_provider_heatmap} to order-sensitive mapping in Sec.~\ref{sec:results_combo}.

\subsection{Confirmatory Test Outcomes (Holm n=30)}

Stage B confirmatory testing yields 5 Holm-significant cells out of 30. Two cells indicate positive improvements for conditional (misinterpretation and rework), while three cells indicate degradation for lazy (update consistency, internal coherence, self-contradiction).

\begin{table}[h]
\centering
\small
\begin{tabular}{llrrl}
\toprule
Pattern & Metric & mean improvement & p(Holm) & Direction \\
\midrule
conditional & misinterpretation\_rate & +0.0781 & 0.0001245 & improved \\
conditional & rework\_cost\_tokens & +122.6481 & 0.0005291 & improved \\
lazy & update\_consistency & -0.1074 & 0.0086253 & degraded \\
lazy & internal\_coherence\_error\_rate & -0.0120 & 0.0095861 & degraded \\
lazy & self\_contradiction\_rate & -0.1062 & 0.0159008 & degraded \\
\bottomrule
\end{tabular}
\caption{Holm-significant Stage B confirmatory cells (Wilcoxon signed-rank, two-sided; family size n=30).}
\label{tab:stageb_holm_sig}
\end{table}

The complete 30-cell test outputs and full six-metric provider-separated means are provided as reproducibility artifacts in the paired CSV files bundled with this manuscript package.

\section{Results II: Order Effects (Combo rep3)}
\label{sec:results_combo}

\subsection{Gate Check}

For the ordered-combination follow-up:
\begin{itemize}
  \item $\tau$ trace pass rate: 0.981481
  \item parse-fail session rate: 0.006944
  \item non-effective rate: 0.716049
\end{itemize}

The follow-up also passes the same gate thresholds.

\subsection{Escalation View: Stable to Unstable Profiles}

To make the boundary contrast explicit, we report providers in escalation order.

\textbf{OpenAI (more order-stable):}
rework improvements stay positive across ordered pairs, in a narrower range (+46.33 to +94.56).
This order stability is not shared by all providers.

\textbf{Anthropic (conditional-gated tendency):}
conditional-first pairs are typically positive on key metrics, while several reverse-order cells are near-zero or slightly negative.

\textbf{Gemini (order-dominated):}
large order-dependent reversals appear in key metrics. For rework cost,
\begin{itemize}
  \item conditional\_branch: +130.11
  \item branch\_conditional: -894.22
\end{itemize}
This yields a swing of +1024.33 tokens from order change alone.

\begin{table}[h]
\centering
\small
\begin{tabular}{lllll}
\toprule
Provider & Pair family & A$\rightarrow$B rework & B$\rightarrow$A rework & Order swing \\
\midrule
OpenAI & conditional/branch & +49.67 & +52.00 & 2.33 \\
Anthropic & conditional/branch & +76.67 & +76.67 & 0.00 \\
Gemini & conditional/branch & +130.11 & -894.22 & 1024.33 \\
\midrule
OpenAI & conditional/hierarchical & +46.56 & +48.22 & 1.67 \\
Anthropic & conditional/hierarchical & +76.67 & -2.67 & 79.33 \\
Gemini & conditional/hierarchical & -845.00 & -79.56 & 765.44 \\
\bottomrule
\end{tabular}
\caption{Ordered-pair rework improvements (tokens). ``Order swing'' is the absolute difference between reversed orders in the same pair family.}
\label{tab:combo_order_rework}
\end{table}

\subsection{Sign-Flip Boundary Map}

From \texttt{processed/all/mst\_sign\_flip\_boundaries.csv}, sign-flip true cells are 14/72.

Notable boundaries:
\begin{itemize}
  \item conditional\_branch: 0/9 flips
  \item hierarchical\_delta: 0/9 flips
  \item branch\_conditional: 5/9 flips
  \item conditional\_hierarchical: 4/9 flips
\end{itemize}

\begin{figure}[h]
\centering
\includegraphics[width=\linewidth]{fig_provider_separated_heatmaps_v9_readable.png}
\caption{Provider-separated heatmaps for ordered combinations (combo rep3, processed/all). Each metric panel uses its own symmetric color scale centered at zero; color intensity is comparable within a panel (not across different metrics). Ordered-pair x-axis labels use short codes (e.g., C->B; C=conditional, H=hierarchical, D=delta, B=branch).}
\label{fig:combo_provider_heatmap}
\end{figure}

\begin{figure}[h]
\centering
\includegraphics[width=\linewidth]{fig_sign_flip_boundaries.png}
\caption{Sign-flip boundary map for ordered combinations (combo rep3, processed/all).}
\label{fig:combo_signflip}
\end{figure}

\subsection{Rep1 to Rep3 Check}

rep3 is consistent with strong provider-by-order interaction, but not every cell preserves rep1 direction. We therefore report this as partial replication with boundary movement, not complete invariance.

\section{Discussion}
\label{sec:discussion}

\subsection{Three-Way Interaction as the Main Result}

The central result is a boundary map, not a universal ranking. Under tested conditions, operator utility depends on provider, pattern, and order jointly. Stage B already shows provider-pattern dependence (Table~\ref{tab:stageb_key}, Fig.~\ref{fig:stageb_provider_heatmap}), while combo rep3 adds order as an additional interaction axis (Table~\ref{tab:combo_order_rework}, Fig.~\ref{fig:combo_provider_heatmap}). The interaction is not subtle in key cells: Gemini changes from +130.11 to -894.22 rework improvement when conditional/branch order is reversed, whereas OpenAI remains positive in both orders with a narrow range.

A practical implication is that operator recommendations without provider/order qualification are incomplete for deployment decisions. ``Use conditional + branch'' is not a sufficient instruction unless order and provider context are explicit.

This also changes how evidence should be summarized. Provider-pooled means are useful as compact diagnostics, but they can hide high-variance directional behavior at the provider level. For deployment governance, the safer reporting primitive is a provider-separated boundary table first, then pooled statistics second. Under this ordering, pooled numbers remain informative while no longer masking where reversals concentrate.

\subsection{Conditional-First Hypothesis (Operational)}

One plausible interpretation is that conditional-first ordering can stabilize the interpretation frame before additional operators increase update complexity. This reading fits the strong separation observed in Gemini for conditional\_branch vs branch\_conditional and the generally positive conditional-first profile in Anthropic.

At the same time, conditional-first is not universally dominant for every cell. In combo rep3, conditional\_hierarchical remains negative on Gemini in both orders for rework, although the magnitude differs by order. This indicates that conditional-leading behavior is a useful tendency in this setup, not a universal safety guarantee.

These are mechanism hypotheses only. The current evidence supports effect observation and boundary mapping, not causal attribution of internal model architecture.

Operationally, this suggests a staged adoption rule rather than a global recommendation: (i) run conditional-first as the default candidate, (ii) verify against the exact provider/order cell before production rollout, and (iii) retain rollback to baseline or single-operator mode when sign-flip risk is detected. This rule is conservative by design and fits the non-replacement stance of NRR: standard prompting remains the fallback, and operator composition is enabled only where the measured cell supports it.

\subsection{Ceiling Effects and Metric Discriminability}

Several confirmatory metrics show many non-effective cells, especially in OpenAI and Anthropic. In this setup, the strongest discriminability is carried by misinterpretation and rework metrics. This does not invalidate the zero-movement metrics; it bounds what can be inferred at c=0.6 under the tested temperatures.

The Holm results are consistent with this pattern. Significant cells concentrate in conditional gains and lazy degradations, while many other cells remain non-significant after correction. This is expected when effect mass is unevenly distributed across metrics and providers. In reporting terms, it argues for showing both inferential outputs (Table~\ref{tab:stageb_holm_sig}) and descriptive boundary outputs (sign-flip maps), instead of relying on a single statistical summary.

From a measurement-design perspective, the implication is to tune tasks for metric discriminability, not only task realism. If a metric family is saturated near zero movement, additional samples mostly tighten confidence around a non-informative region. Under this setup, expanding task complexity and scenario diversity is likely to improve inferential value more than simply increasing repeats in the same low-variance cells.

\subsection{Unmapped Regions and Extension Scope}

The current map remains partial. Unmeasured regions include additional temperatures for ordered combinations, task families beyond spec-to-code, and longer operator chains. We therefore treat this manuscript as a bounded map of one tested region, not an exhaustive catalog.

Future extension should preserve the same paired protocol and boundary-first reporting policy, so newly mapped cells remain directly comparable to the present results. Priority extensions are: (i) combo runs at t=0.3, (ii) task-family transfer checks, and (iii) three-operator chains where order permutations expand rapidly. These are natural next steps for the same map-building program rather than separate methodological proposals.

To keep claim strength aligned with evidence growth, each extension should predeclare whether it is confirmatory or exploratory before execution. In practice, this means fixing families and multiplicity rules for confirmatory expansions, while keeping newly opened regions descriptive until they are independently replicated. This separation preserves interpretability as the map scales.

\section{Limitations}
\label{sec:limitations}

\begin{itemize}
  \item Fixed task family (spec-to-code) and fixed complexity setting (c=0.6).
  \item Ordered follow-up uses t=0.0 only.
  \item Provider versions are time-bound and may shift with future model updates.
  \item Several metrics show high non-effective regions, limiting discriminability.
\end{itemize}

\section{Conclusion}
\label{sec:conclusion}

Under tested conditions, NRR operator effects are condition-dependent rather than intrinsic constants. The same operators can improve or degrade depending on provider and order. In the ordered-combination follow-up, conditional\_branch showed 0/9 sign-flip metrics while branch\_conditional showed 5/9, illustrating how reversal risk can change with order alone. The contribution of this paper is an explicit boundary map that makes those dependencies auditable and reproducible.

\section{Reproducibility}
\label{sec:repro}

\begin{itemize}
  \item Stage B run: \texttt{results/main\_c06\_5patterns\_3providers\_20260302T054857Z}
  \item Combo rep3 run: \texttt{results/combo\_probe\_c06\_ordered8\_rep3\_20260303T020142Z}
  \item Main scripts: \texttt{scripts/run\_mst\_c06\_main\_5patterns\_3providers.sh}, \texttt{scripts/run\_mst\_c06\_combo\_probe\_ordered8\_rep3.sh}
  \item Confirmatory test script: \texttt{repro/compute\_stageb\_confirmatory\_tests.py}
\end{itemize}

\section*{Acknowledgments}
The author gratefully acknowledges assistance from Claude (Anthropic), ChatGPT/Codex (OpenAI), and Gemini (Google) for language refinement and drafting support. All claims and decisions in this manuscript remain the sole responsibility of the author.

\begin{thebibliography}{99}

\bibitem{saito2025nrr}
K.~Saito,
\newblock ``Non-Resolution Reasoning (NRR): A Structural Framework for Deferred Semantic Fixation,''
\newblock \emph{arXiv preprint arXiv:2512.13478}, 2025.

\bibitem{saito2026phi}
K.~Saito,
\newblock ``NRR-Phi: Positional Semantic Dynamics Through Non-Resolution Operators,''
\newblock \emph{arXiv preprint arXiv:2601.19933}, 2026.

\bibitem{saito2026ime}
K.~Saito,
\newblock ``NRR-IME: Implementation Patterns for Stateless API Systems,''
\newblock Unpublished manuscript (submitted, under moderation), 2026.

\bibitem{liang2022helm}
P.~Liang et al.,
\newblock ``Holistic Evaluation of Language Models,''
\newblock \emph{Transactions on Machine Learning Research}, 2023.

\bibitem{zheng2023mtbench}
L.~Zheng et al.,
\newblock ``Judging LLM-as-a-Judge with MT-Bench and Chatbot Arena,''
\newblock \emph{NeurIPS Datasets and Benchmarks}, 2023.

\bibitem{dubois2024alpacaeval}
Y.~Dubois et al.,
\newblock ``AlpacaEval 2.0: Evaluating Language Models as Distributional Samplers,''
\newblock \emph{arXiv preprint arXiv:2404.04475}, 2024.

\bibitem{wei2022cot}
J.~Wei et al.,
\newblock ``Chain-of-Thought Prompting Elicits Reasoning in Large Language Models,''
\newblock \emph{NeurIPS}, 2022.

\bibitem{kojima2022zeroshot}
T.~Kojima et al.,
\newblock ``Large Language Models are Zero-Shot Reasoners,''
\newblock \emph{NeurIPS}, 2022.

\bibitem{wang2023selfconsistency}
X.~Wang et al.,
\newblock ``Self-Consistency Improves Chain of Thought Reasoning in Language Models,''
\newblock \emph{ICLR}, 2023.

\bibitem{zhao2021calibrate}
Z.~Zhao et al.,
\newblock ``Calibrate Before Use: Improving Few-Shot Performance of Language Models,''
\newblock \emph{ICML}, 2021.

\end{thebibliography}

\end{document}
